% !TEX root = ../main.tex
%% Detection Rules and Parser
\section{Detection Rules and Parser}
\label{app:rules}

Detection was deterministic. This appendix records the rules that raised the alerts
used in the experiment, together with the mapping from Teleport audit events to the
fields the rules read. Detection is not an evaluated claim of this thesis; the rules
are documented so that the alert set is reproducible.

\subsection{Teleport audit event to alert field mapping}
\label{app:rules:mapping}

Teleport Community Edition writes one JSON object per audit event. Five event codes
are relevant to database access, and two carry SQL text. The lab corpus was generated
through \code{psql}, which uses the simple query protocol, so every statement arrived
as \code{TDB02I}. In the production tenant described in
Section~\ref{sec:method:secops}, 21 of 25 sampled statement events were instead
\code{TPG00I}, because application connection pools use prepared statements. A
deployment of these rules must therefore match both event types; the parse event
carries the statement text with \code{\$n} placeholders in place of literal values,
which leaves every predicate below evaluable except those that test a literal.

\begin{table}[H]
\centering
\caption{Teleport database audit event codes}
\label{tab:app:eventcodes}
\small
\begin{tabular}{@{}llp{6.2cm}@{}}
\toprule
Code & Event & Meaning \\
\midrule
\code{TDB00I} & \code{db.session.start}       & Session to the database established \\
\code{TDB00W} & \code{db.session.start}       & Session denied by proxy authorisation \\
\code{TDB01I} & \code{db.session.end}         & Session closed \\
\code{TDB02I} & \code{db.session.query}       & A single SQL statement was proxied (simple query protocol) \\
\code{TPG00I} & \code{db.session.postgres.statements.parse} & A prepared statement was parsed (extended query protocol); production only, not exercised by the lab corpus \\
\bottomrule
\end{tabular}
\end{table}

Table~\ref{tab:app:udmmap} gives the mapping the rules read. It was verified against
the production tenant rather than taken from documentation: a sample of 25 live
events under the \code{TELEPORT\_ACCESS\_PLANE} log type was read through the search
API, only field paths and non-identifying metadata were retained, and no event
content was stored. The log type is the native one, but the tenant's
\code{parserVersion} reads \code{custom-parser}, so the field placement below is that
of the organisation's parser override and not of the stock parser. The first draft of
this appendix placed the statement text in a \code{db\_query} resource label; no
such label exists in the live events, and the rule sources were rewritten to the
verified paths. The originals are retained beside them for traceability. Because the
experiment was executed against the local corpus rather than through ingestion, the
same predicates are transcribed line for line into the Python evaluator, and the
YARA-L sources below are the platform-native expression of what was actually run.

\begin{table}[H]
\centering
\caption{Teleport audit field to UDM field mapping}
\label{tab:app:udmmap}
\small
\begin{tabular}{@{}llp{4.1cm}@{}}
\toprule
Teleport field & UDM field (verified) & Purpose in rules \\
\midrule
\code{event}          & \code{metadata.product\_event\_type}             & Source scoping \\
\code{code}           & \code{additional.fields["code"]}                 & Event code \\
\code{user}           & \code{principal.user.userid}                     & The named human \\
\code{db\_user}       & \code{target.user.user\_display\_name}           & Database role assumed \\
\code{db\_name}       & \code{...labels["db\_name"]}                     & Database \\
\code{db\_service}    & \code{target.application}                        & Proxied service \\
\code{db\_query}      & \code{security\_result.description}              & Statement text \\
\code{sid}            & \code{network.session\_id}                       & Session correlation \\
\code{time}           & \code{metadata.event\_timestamp}                 & Hour-of-day tests \\
(none)                 & \code{additional.fields["rows\_affected"]}       & Row count; database-side enrichment, absent from Teleport \\
\bottomrule
\end{tabular}
\end{table}

Two properties of this mapping constrain what the rules can express. First,
Teleport's \code{success} records whether the proxy's own role-based access control
permitted the statement, not whether PostgreSQL executed it; a statement the proxy
allowed and the database refused is logged as successful. Ground truth therefore
records the database-side outcome separately. Second, the stream carries statements
and not result sets, so the number of rows a statement returned or changed is not
observable from Teleport at all. Four rules (\code{PG-EXF-001}, \code{PG-EXF-003},
\code{PG-TAMP-001} and \code{PG-TAMP-004}) nevertheless test a row count, and the
alert evidence shown to the remediation arms carries one. In the lab that value came
from the traffic generator's own record, joined to the audit event by the statement
marker, standing in for the database-side enrichment a deployment would have to
build. To measure how much the detection layer leaned on it, the evaluator was re-run
with the row count hidden: step-level recall fell from 35 to 25 of 60, session-level
recall from 12 to 10 of 12, and benign alerting from six sessions to three. The two
attack sessions lost are the vendor banking sweep (\code{PG-EXF-001}) and the sealed-bid
tampering (\code{PG-TAMP-004}); three of the six false-positive controls also stop
alerting. A proxy-only deployment of these rules would therefore have produced thirteen
alerts rather than eighteen, and would have lost half of the benign alerts on which the
generative arm's restraint was measured. Section~\ref{sec:discussion} carries this as a limit on how far
the alert set represents what \secops{} alone would raise.

\subsection{Rules}
\label{app:rules:list}

Eleven rules were written, covering the three malicious behaviour classes. Ten are
single-event; \code{PG-EXF-003} is session-scoped and aggregates over a one-hour match
window. Table~\ref{tab:app:ruleindex} indexes them; the full YARA-L~2.0 sources follow.

\begin{table}[H]
\centering
\caption{Detection rule index}
\label{tab:app:ruleindex}
\small
\footnotesize
\setlength{\tabcolsep}{4pt}
\begin{tabular}{@{}lp{4.4cm}cllrr@{}}
\toprule
Rule & Intent & Class & Severity & ATT\&CK & Alerts & TP \\
\midrule
PG-EXF-001   & Unbounded read of sensitive columns    & S1 & High     & T1213 & 2 & 1 \\
PG-EXF-002   & Bulk export via \code{COPY TO STDOUT}  & S1 & High     & T1213, T1005 & 3 & 2 \\
PG-EXF-003   & Out-of-hours paginated bulk read       & S1 & High     & T1213, T1030 & 2 & 1 \\
PG-PRIV-001  & Grant to self or lateral role grant    & S2 & Critical & T1078.004 & 2 & 1 \\
PG-PRIV-002  & Role altered to \code{SUPERUSER}       & S2 & Critical & T1098 & 1 & 1 \\
PG-PRIV-003  & Shadow role or default privilege change& S2 & High     & T1136.001 & 3 & 2 \\
PG-TAMP-001  & \code{DELETE}/\code{UPDATE} unbounded  & S3 & Critical & T1565.001 & 1 & 0 \\
PG-TAMP-002  & \code{TRUNCATE} or \code{DROP}         & S3 & Critical & T1485 & 1 & 1 \\
PG-TAMP-003  & Obfuscated SQL construction            & S3 & High     & T1027 & 1 & 1 \\
PG-TAMP-004  & Mass update of sealed bid amounts      & S3 & Critical & T1565.001 & 1 & 1 \\
PG-ANTIF-001 & Audit log tampering                    & S3 & Critical & T1070 & 1 & 1 \\
\midrule
\multicolumn{5}{@{}l}{Total} & 18 & 12 \\
\bottomrule
\end{tabular}
\end{table}

The \emph{Alerts} column gives the number of alerts each rule raised over the whole
corpus and \emph{TP} the number of those that were true positives, that is, that fired
on a session belonging to class S1, S2 or S3. The six alerts that are not true
positives all fired on class B1 sessions, which is what B1 exists to produce. No rule
fired on any of the \num{588} class B0 sessions. The six false positives are spread
one apiece across six different rules (\code{PG-EXF-001}, \code{PG-EXF-002},
\code{PG-EXF-003}, \code{PG-PRIV-001}, \code{PG-PRIV-003} and \code{PG-TAMP-001}),
which is the pattern to expect when the benign-but-alerting sessions were designed to
exercise a different detector each. \code{PG-TAMP-001} is the only rule whose every
alert was a false positive: its single firing was the scheduled retention job,
\code{DELETE FROM payments WHERE paid\_at < now() - interval '700 days'}, which
removed \num{1263} rows. The statement is predicated and entirely routine, but the
rule keys on rows affected rather than on the presence of a \code{WHERE} clause, and
at that volume a legitimate retention sweep and a sabotage attempt are the same event
at the statement level. Separating them requires the change window, which is supplied
to the remediation stage as \code{change\_window} but is not available to the rule.

\subsection{Correspondence with the production rule set}
\label{app:rules:production}

The production tenant carries fifteen Teleport rules written by the organisation's
security team, twelve of them for database sessions, all but one enabled. They were
exported read-only. Table~\ref{tab:app:correspondence} sets each lab rule against its
nearest deployed counterpart and quotes the deployed predicate in abridged form. The
pattern is consistent: the deployed rules are generic, multi-engine keyword matchers
over \code{security\_result.description}, while the lab rules are PostgreSQL-specific,
role-aware and grounded in the schema under protection. Three deployed database rules
have no lab counterpart because they address concerns outside the insider classes:
first-time access, sessions without multi-factor authentication, and lateral movement
across databases. A fourth, \code{teleport\_db\_sql\_injection\_probe}, matches
\code{UNION SELECT} with a comment terminator, boolean and timing probes and hex
encoding. It is the deployed trace of the framing the research proposal used, and it
is deliberately not reproduced in the lab for the reason given in
Appendix~\ref{app:scenarios}: those patterns belong to an unauthenticated attacker
reaching the database through an application, not to an authorised principal on a
proxied session. Table~\ref{tab:app:proddetections} gives the detections each deployed
rule raised over the thirty days to 7~September~2026, as a real-world anchor for alert
volume; only counts were read.

\begin{table}[H]
\centering
\caption{Lab rules against their nearest deployed counterparts}
\label{tab:app:correspondence}
\scriptsize
\setlength{\tabcolsep}{4pt}
\providecommand{\pcode}[1]{{\ttfamily\hspace{0pt}#1}}
\begin{tabular}{@{}>{\raggedright\arraybackslash}p{1.6cm}>{\raggedright\arraybackslash}p{3.2cm}>{\raggedright\arraybackslash}p{3.7cm}>{\raggedright\arraybackslash}p{5.5cm}@{}}
\toprule
Lab rule & Deployed counterpart & Deployed predicate (abridged) & Difference \\
\midrule
PG-EXF-001 & \pcode{teleport\_\allowbreak{}db\_\allowbreak{}sensitive\_\allowbreak{}data} & regex \pcode{credentials|\allowbreak{}api\_\allowbreak{}key|\allowbreak{}credit\_\allowbreak{}card|\allowbreak{}secret} & Deployed keys on generic column names, none of which exist in this schema; lab names the schema's sensitive columns and requires no \pcode{LIMIT} \\
PG-EXF-002 & none & & \pcode{COPY \dots{} TO STDOUT} is uncovered; \pcode{teleport\_\allowbreak{}db\_\allowbreak{}file\_\allowbreak{}os\_\allowbreak{}ops} matches only \pcode{COPY \dots{} PROGRAM} \\
PG-EXF-003 & \pcode{teleport\_\allowbreak{}db\_\allowbreak{}data\_\allowbreak{}exfiltration}; \pcode{teleport\_\allowbreak{}db\_\allowbreak{}after\_\allowbreak{}hours} & more than 200 events per user in 5 minutes; 17:00--22:00 UTC window & Deployed exfiltration is a query-count threshold with no notion of sensitivity; S1 sessions issue three to seven statements and would not trip it \\
PG-PRIV-001 & \pcode{teleport\_\allowbreak{}db\_\allowbreak{}privilege\_\allowbreak{}escalation} & regex \pcode{GRANT (ALL|\allowbreak{}SELECT|\allowbreak{}\dots)} & Deployed fires on every \pcode{GRANT} by any role; lab restricts to non-administrative roles and self-grants \\
PG-PRIV-002 & \pcode{teleport\_\allowbreak{}db\_\allowbreak{}privilege\_\allowbreak{}escalation} & regex \pcode{ALTER (USER|\allowbreak{}ROLE)} & Deployed fires on any role alteration; lab requires a privilege attribute such as \pcode{SUPERUSER} \\
PG-PRIV-003 & \pcode{teleport\_\allowbreak{}db\_\allowbreak{}privilege\_\allowbreak{}escalation} & regex \pcode{CREATE (USER|\allowbreak{}ROLE)} & \pcode{ALTER DEFAULT PRIVILEGES} is uncovered in the deployed set \\
PG-TAMP-001 & \pcode{teleport\_\allowbreak{}db\_\allowbreak{}destructive\_\allowbreak{}query\_\allowbreak{}high}; \pcode{teleport\_\allowbreak{}db\_\allowbreak{}critical\_\allowbreak{}write} & regex \pcode{DELETE FROM}; \pcode{UPDATE \dots{} SET} on named production databases & Deployed has no \pcode{WHERE}-clause or volume test, so every delete and update on a production database alerts \\
PG-TAMP-002 & \pcode{teleport\_\allowbreak{}db\_\allowbreak{}destructive\_\allowbreak{}query\_\allowbreak{}critical} & regex \pcode{DROP TABLE|\allowbreak{}TRUNCATE TABLE|\allowbreak{}ALTER TABLE} & Closest match in the set; deployed also fires on schema migration \\
PG-TAMP-003 & none & & Comment-split keywords and \pcode{chr()} construction are uncovered; case mangling is caught incidentally by \pcode{nocase} \\
PG-TAMP-004 & \pcode{teleport\_\allowbreak{}db\_\allowbreak{}critical\_\allowbreak{}write} & regex \pcode{UPDATE \dots{} SET} on named databases & Deployed fires on every update; lab requires the \pcode{bids} table, a sealed-value column and volume \\
PG-ANTIF-001 & \pcode{teleport\_\allowbreak{}db\_\allowbreak{}destructive\_\allowbreak{}query\_\allowbreak{}high} & regex \pcode{DELETE FROM} & Deployed is not audit-specific; \pcode{UPDATE} of \pcode{audit\_\allowbreak{}log} rows is covered only by the generic write rule \\
\bottomrule
\end{tabular}
\end{table}

% !TEX root = ../main.tex
%% generated from a 30-day detection count export of the deployed tenant (counts only; no detection content was read)
\begin{table*}[t]
\centering
\caption{Detections raised by the deployed Teleport database rules, 8 August to 7 September 2026}
\label{tab:app:proddetections}
\small
\begin{tabular}{@{}lccr@{}}
\toprule
Deployed rule & Enabled & Cadence & Detections \\
\midrule
\code{teleport\_db\_after\_hours} & yes & hourly & $\geq$\num{25000} \\
\code{teleport\_db\_critical\_write} & yes & hourly & \num{18906} \\
\code{teleport\_db\_data\_exfiltration} & yes & hourly & \num{1228} \\
\code{teleport\_db\_destructive\_query\_critical} & yes & hourly & \num{216} \\
\code{teleport\_db\_destructive\_query\_high} & yes & hourly & \num{4493} \\
\code{teleport\_db\_file\_os\_ops} & yes & live & \num{0} \\
\code{teleport\_db\_first\_time\_access} & no & hourly & \num{0} \\
\code{teleport\_db\_lateral\_movement} & yes & hourly & \num{143} \\
\code{teleport\_db\_privilege\_escalation} & yes & live & \num{187} \\
\code{teleport\_db\_sensitive\_data} & yes & hourly & \num{0} \\
\code{teleport\_db\_sql\_injection\_probe} & yes & live & \num{0} \\
\code{teleport\_db\_without\_mfa} & yes & hourly & \num{14931} \\
\midrule
\multicolumn{3}{@{}l}{Total, twelve rules, thirty days} & $\geq$\num{65104} \\
\multicolumn{3}{@{}l}{Lab, eleven rules, one simulated month} & \num{18} \\
\bottomrule
\end{tabular}
\begin{minipage}{0.92\linewidth}
\vspace{2pt}\scriptsize
Counts were paged from the detection search API in pages of \num{1000} and capped at
25 pages, so the after-hours figure is a floor. Only counts were retained. The
tenant's SOAR case search returned no cases over the same period, so none of these
detections entered a case workflow with a human approval step.
\end{minipage}
\end{table*}


\subsection{Rule sources}
\label{app:rules:sources}

The sources are reproduced verbatim from \code{lab/rules/}, in the field mapping of
Table~\ref{tab:app:udmmap}. The first-draft sources that used the unverified
\code{db\_query} label are kept in \code{lab/rules/\_v1\_labels/} and differ only in
field paths.

\subsubsection{Data harvesting (S1)}
\lstinputlisting[language=yaral,caption={PG-EXF-001}]{../testbed/rules/pg-exf-001.yaral}
\lstinputlisting[language=yaral,caption={PG-EXF-002}]{../testbed/rules/pg-exf-002.yaral}
\lstinputlisting[language=yaral,caption={PG-EXF-003}]{../testbed/rules/pg-exf-003.yaral}

\subsubsection{Privilege escalation (S2)}
\lstinputlisting[language=yaral,caption={PG-PRIV-001}]{../testbed/rules/pg-priv-001.yaral}
\lstinputlisting[language=yaral,caption={PG-PRIV-002}]{../testbed/rules/pg-priv-002.yaral}
\lstinputlisting[language=yaral,caption={PG-PRIV-003}]{../testbed/rules/pg-priv-003.yaral}

\subsubsection{Destructive or obfuscated SQL (S3)}
\lstinputlisting[language=yaral,caption={PG-TAMP-001}]{../testbed/rules/pg-tamp-001.yaral}
\lstinputlisting[language=yaral,caption={PG-TAMP-002}]{../testbed/rules/pg-tamp-002.yaral}
\lstinputlisting[language=yaral,caption={PG-TAMP-003}]{../testbed/rules/pg-tamp-003.yaral}
\lstinputlisting[language=yaral,caption={PG-TAMP-004}]{../testbed/rules/pg-tamp-004.yaral}
\lstinputlisting[language=yaral,caption={PG-ANTIF-001}]{../testbed/rules/pg-antif-001.yaral}

\subsection{Rule hit rate over the injected attacks}
\label{app:rules:hitrate}

% !TEX root = ../main.tex
%% generated from data/corpus/detection_hitrate.csv
\begin{table}[t]
\centering
\caption{Rule hit rate over the injected attack steps, and alerting on benign activity}
\label{tab:detection_hitrate}
\footnotesize
\setlength{\tabcolsep}{4pt}
\begin{tabular}{@{}lrrrrr@{}}
\toprule
\textbf{Class} & \textbf{Steps} & \textbf{Alerted} & \textbf{Missed} & \textbf{Hit rate} & \textbf{Benign} \\
\midrule
S1 Harvesting          & 20 & 10 & 10 & 50.0\% & 3 \\
S2 Priv.\ escalation   & 20 & 11 &  9 & 55.0\% & 2 \\
S3 Destructive         & 20 & 14 &  6 & 70.0\% & 1 \\
\midrule
Total                         & 60 & 35 & 25 & 58.3\% & 6 \\
\bottomrule
\end{tabular}

\vspace{2pt}
\scriptsize\raggedright Session-level detection was 12 of 12: every attack session raised at
least one alert. All six benign alerts came from the B1 class, which was designed to
trip a detector. None of the 588 B0 routine sessions raised an alert.
\end{table}


Every attack session was detected: session-level recall was 12 of 12. At the level of
individual statements, however, just over four steps in ten went undetected. That is the expected consequence
of writing rules that are specific rather than sensitive: every rule fires on a single
statement or a bounded session window, so reconnaissance and verification statements such as
\code{SELECT current\_user} and \code{SELECT count(*)} pass without comment. Those
are in fact all of the missed steps, and the rules are arguably right to ignore them,
since the same statements are issued constantly by legitimate users. The design accepts this, because the object of study is what happens after an
alert is raised, and a rule set tuned for recall would have flooded the queue with
routine traffic and changed the base rate the corpus was built to hold. The
complementary result is that none of the 588 routine sessions alerted; all six benign
alerts came from the B1 class, which exists precisely to produce them.
